% Vertiefte Lesefassung. Die nummerierten Aussagen bleiben im Hauptband.
\hypertarget{index-period.reading}{}
\section*{Was eine endliche Potenzfolge verrät}

Wir wählen ein Element \(a\) einer endlichen Halbgruppe \((A,\star)\)
und multiplizieren es immer wieder mit sich selbst:
\[
 a,\quad a^2,\quad a^3,\quad a^4,\quad\ldots
\]
Die Exponenten wachsen unbegrenzt, doch die Werte liegen in der endlichen
Menge \(A\). Deshalb muss irgendwann ein schon bekannter Wert wiederkehren.
Entscheidend ist, dass sich von diesem Augenblick an auch der weitere Verlauf
wiederholt. Die Folge besteht aus einem endlichen Anlauf und einem Kreis.

Diese Lesefassung erklärt, wie man beide Teile eindeutig bestimmt,
wie sie die von \(a\) erzeugte Unterhalbgruppe beschreiben und warum
unter den Potenzen stets ein idempotentes Element vorkommt.
Definitionen und nummerierte Sätze stehen weiterhin in Band~37.
Die \IndexPeriodProofsLink{} enthalten die genauen Aussagen mit ihren
formalen Ableitungen; die \IndexPeriodMainLink{} führt zum Hauptband.

Eine Halbgruppe besitzt eine abgeschlossene, assoziative Multiplikation.
Ein neutrales Element wird nicht vorausgesetzt. Daher beginnen wir
bei den \emph{positiven} Potenzen:
\[
 a^1=a,\qquad a^{r+1}=a^r\star a\quad(r\geq1).
\]
Aus der Assoziativität folgt das Additionsgesetz
\[
 a^{r+s}=a^r\star a^s\qquad(r,s\geq1).
\]
Diese Grundlagen werden in Band~28 bewiesen. Wir verwenden keine
Potenz \(a^0\). Ein Folgenindex \(0\) kann zwar den ersten Eintrag
bezeichnen; dieser Eintrag ist dann \(a^1\), nicht \(a^0\).

\section*{Ein Beispiel mit Anlauf und Kreis}

Seien \(a,b,c\) drei verschiedene Elemente. Wir betrachten dieselbe
Halbgruppe wie in der Lesefassung zu Multiplikationstafeln:
\[
 \begin{array}{c|ccc}
 \star&a&b&c\\\hline
 a&b&c&b\\
 b&c&b&c\\
 c&b&c&b
 \end{array}.
\]
Die Assoziativität lässt sich an einer gemeinsamen Regel ablesen.
Setze \(q(a)=q(c)=1\) und \(q(b)=0\). Jedes Produkt liegt in \(\{b,c\}\);
es ist genau dann \(b\), wenn die Summe der beiden \(q\)-Werte gerade ist.
Damit hat jede Klammerung eines Dreifachprodukts den \(q\)-Wert
\(q(x)+q(y)+q(z)\) modulo \(2\). Auf \(\{b,c\}\) unterscheidet \(q\)
die Elemente, also stimmen die beiden Klammerungen überein.

Die Potenzfolge von \(a\) lautet:
\[
 \begin{array}{c|cccccccc}
 r&1&2&3&4&5&6&7&8\\\hline
 a^r&a&b&c&b&c&b&c&b
 \end{array}.
\]
Vor dem Kreis steht genau ein Wert, nämlich \(a\). Danach wechseln
sich \(b\) und \(c\) ab. Der Kreis beginnt beim Exponenten \(2\)
und besitzt die Länge \(2\). Wir werden dafür
\[
 \mu=2,\qquad\lambda=2
\]
schreiben. Die Zahl \(\mu\) bezeichnet die \emph{erste Stelle im Kreis}.
Die Länge des davorliegenden Anlaufs beträgt deshalb \(\mu-1\).

\section*{Die erste Wiederholung bestimmt die Parameter}

Hat \(A\) genau \(N\) Elemente, können die \(N+1\) Potenzwerte
\(a^1,\ldots,a^{N+1}\) nicht alle verschieden sein. Es gibt
positive Exponenten \(i<j\) mit \(a^i=a^j\).
Die im Hauptband definierte Menge
\[
 \mathcal R_{A,\star}(a)
 =\{j\geq1\mid \text{für ein }i\text{ mit }1\leq i<j
                    \text{ gilt }a^i=a^j\}
\]
sammelt die \emph{Endpunkte} solcher Wiederholungen.
Sie ist nichtleer; dies ist
\FormulaRefAuto{FiniteSemigroupPowerRepetitionExists}[thm].
\IndexPeriodProofLink{FiniteSemigroupPowerRepetitionExists}.

Wir wählen das kleinste \(j\) in dieser Menge.
Dann sind \(a^1,\ldots,a^{j-1}\) paarweise verschieden.
Andernfalls wäre schon ein früherer Exponent ein Wiederholungsendpunkt.
Zu \(a^j\) gibt es unter diesen Werten daher genau einen Treffer
\(a^i\). Nun setzen wir
\[
 \mu=i,\qquad\lambda=j-i.
\]
Beide Zahlen sind positiv, und es gilt
\[
 a^{\mu+\lambda}=a^\mu.
\]
Vor diesem Endpunkt sind noch alle Potenzwerte verschieden:
\[
 1\leq r<s<\mu+\lambda\quad\Longrightarrow\quad a^r\ne a^s.
\]
Positivität, Wiederholungsgleichung und diese Trennung bilden gemeinsam
den Begriff des Index--Perioden-Paars
\FormulaRefAuto{FiniteSemigroupIndexPeriodPairDef}[def].

Die Trennungsbedingung gehört zur Definition dazu. Im Dreierbeispiel
gilt zwar auch \(a^2=a^6\); das Paar \((2,4)\) beschreibt aber nicht
den ersten Kreis. Schon bei Exponent \(4\) war \(b\) erneut aufgetreten.
Das wirkliche Paar ist \((2,2)\), und der erste Wiederholungsendpunkt
ist \(\mu+\lambda=4\).

\section*{Warum Index und Periode eindeutig sind}

Der Index--Perioden-Satz besagt: Zu jedem Element einer endlichen
Halbgruppe gibt es genau ein solches Paar
\FormulaRefAuto{FiniteSemigroupIndexPeriod}[thm].
Die Existenz folgt aus der eben beschriebenen ersten Wiederholung.
Für die Eindeutigkeit betrachten wir zwei mögliche Paare
\((\mu,\lambda)\) und \((\mu',\lambda')\).

Zuerst müssen ihre Endpunkte gleich sein. Wäre etwa
\(\mu+\lambda<\mu'+\lambda'\), dann lägen die beiden verschiedenen
Exponenten \(\mu\) und \(\mu+\lambda\) noch vor dem zweiten Endpunkt.
Ihre Potenzen sind nach der ersten Wiederholungsgleichung gleich.
Das widerspricht der Trennungsbedingung des zweiten Paars.
Die umgekehrte Ungleichung ist ebenso ausgeschlossen.

Sei also \(E=\mu+\lambda=\mu'+\lambda'\).
Dann gilt
\[
 a^\mu=a^E=a^{\mu'}.
\]
Beide Exponenten liegen zwischen \(1\) und \(E-1\).
Da die Potenzen dort paarweise verschieden sind, folgt \(\mu=\mu'\).
Aus den gleichen Summen folgt schließlich \(\lambda=\lambda'\).

Die Beweistabellen trennen diese Überlegung in die Eindeutigkeit
des Endpunkts, die Eindeutigkeit des Index und die Eindeutigkeit
des gesamten Paars.
\IndexPeriodProofLink{FiniteSemigroupIndexPeriod}.

\section*{Warum der Kreis unbegrenzt weiterläuft}

Aus einer Potenzgleichheit
\(a^{\mu+\lambda}=a^\mu\) folgt für jeden positiven Exponenten \(t\):
\[
 a^{(\mu+\lambda)+t}
 =a^{\mu+\lambda}\star a^t
 =a^\mu\star a^t
 =a^{\mu+t}.
\]
Für \(t=0\) ist die Aussage bereits die Ausgangsgleichheit.
Dieser Fall benötigt keine Multiplikation mit \(a^0\).

Jeder Exponent \(r\geq\mu\) lässt sich als \(\mu+t\) mit \(t\geq0\)
schreiben. Somit gilt
\[
 \boxed{a^{r+\lambda}=a^r\qquad(r\geq\mu).}
\]
Das ist die periodische Fortsetzung
\FormulaRefAuto{FiniteSemigroupPowerEventuallyPeriodic}[thm].
Durch wiederholtes Anwenden folgt auch
\[
 a^{r+q\lambda}=a^r\qquad(r\geq\mu,\ q\geq0),
\]
wie in \FormulaRefAuto{FiniteSemigroupMultiplePeriodShift}[thm]
formal bewiesen wird.
\IndexPeriodProofLink{FiniteSemigroupMultiplePeriodShift}.

Im Dreierbeispiel ist \(a^{r+2}=a^r\) für alle \(r\geq2\).
Die Einschränkung \(r\geq\mu\) ist wesentlich:
Bei \(r=1\) wäre \(a^3=a\) falsch, denn \(a^3=c\).
Die Wiederholung setzt erst beim Eintritt in den Kreis ein.

\section*{Alle Potenzen auf einen endlichen Abschnitt zurückführen}

Jede positive Potenz stimmt mit einem der Werte
\[
 a^1,\ldots,a^{\mu+\lambda-1}
\]
überein. Für kleinere Exponenten ist nichts zu tun.
Ist \(r\geq\mu+\lambda\), so darf man \(\lambda\) abziehen:
Der neue Exponent liegt noch mindestens bei \(\mu\), und die
periodische Fortsetzung liefert
\[
 a^r=a^{r-\lambda}.
\]
Solange der Exponent mindestens \(\mu+\lambda\) beträgt, wiederholt
man diesen Schritt. Er endet, weil der positive Exponent jeweils
kleiner wird. Für einen Exponenten aus dem Kreis erhält man
schließlich genau einen Vertreter zwischen \(\mu\) und
\(\mu+\lambda-1\).

Die von \(a\) erzeugte Unterhalbgruppe besteht aus allen positiven
Potenzen von \(a\). Daher gilt
\[
 \langle a\rangle
 =\{a^1,\ldots,a^{\mu+\lambda-1}\},
 \qquad
 \boxed{|\langle a\rangle|=\mu+\lambda-1.}
\]
Die Anzahl folgt nicht allein aus der Reduktion: Zusätzlich sind
die aufgeführten Werte nach der Trennungsbedingung verschieden.
Beide Teile werden in
\FormulaRefAuto{FiniteSemigroupMonogenicSize}[thm] zusammengeführt.
\IndexPeriodProofLink{FiniteSemigroupMonogenicSize}.

Damit lassen sich auch alle Produkte innerhalb dieser Unterhalbgruppe
berechnen. Man addiert die Exponenten und reduziert anschließend:
\[
 a^r\star a^s=a^{r+s}.
\]
Im Beispiel ist \(b=a^2\) und \(c=a^3\). Deshalb
\[
 b\star c=a^5=a^3=c,\qquad c\star c=a^6=a^2=b.
\]
Hier ist \(\langle a\rangle=A\).
In einer anderen Halbgruppe kann die erzeugte Unterhalbgruppe kleiner
sein. Dann bestimmen Index und Periode dieses einen Elements nur
die Produkte seiner Potenzen.

\section*{Eine idempotente Potenz finden}

Ein Element \(e\) heißt \emph{idempotent}, wenn \(e\star e=e\).
Wir suchen eine Potenz \(e=a^n\) mit
\[
 a^{2n}=a^n.
\]
Dazu wählen wir einen positiven Exponenten \(n\), der mindestens
\(\mu\) groß und ein Vielfaches von \(\lambda\) ist.
Eine stets mögliche Wahl lautet
\[
 n=\lambda\mu.
\]
Weil \(\lambda\geq1\) und \(\mu\geq1\), ist \(n\) positiv
und \(n\geq\mu\). Weil \(n\) ein Periodenvielfaches ist, dürfen
wir vom Exponenten \(n\) aus um weitere \(n\) Stellen verschieben:
\[
 (a^n)^2=a^{2n}=a^n.
\]
Damit besitzt jedes Element einer endlichen Halbgruppe eine
idempotente Potenz
\FormulaRefAuto{FiniteSemigroupIdempotentPower}[thm].
\IndexPeriodProofLink{FiniteSemigroupIdempotentPower}.

Für das Dreierbeispiel genügt bereits \(n=2\):
Es ist \(2\geq\mu\), und \(2\) ist ein Vielfaches von \(\lambda\).
Wir finden \(a^2=b\) mit \(b^2=b\).
Die allgemeine Wahl \(n=\lambda\mu=4\) liefert denselben Wert
\(a^4=b\). Sie ist eine sichere Wahl, muss aber nicht die
kleinste mögliche sein.

Ist die Halbgruppe nichtleer, wählen wir ein Element und erhalten
auf diesem Weg wenigstens ein Idempotentes. Das ist
\FormulaRefAuto{FiniteSemigroupIdempotentExists}[thm].
Beim leeren Träger ist diese Wahl unmöglich; deshalb gehört die
Voraussetzung \(A\ne\varnothing\) ausdrücklich zu diesem Satz.

Ein Idempotentes ist nicht automatisch ein neutrales Element der
ganzen Halbgruppe. Im Beispiel gilt zwar \(b^2=b\), aber
\(b\star a=c\ne a\). Die Halbgruppe hat somit ein Idempotentes,
ohne dass dieses alle Elemente unverändert lässt.

\section*{Drei Randfälle zum Einordnen}

\begin{description}
 \item[Kein Anlauf: \(\mu=1\).]
 Die erste Potenz liegt bereits im Kreis. In der additiv
 geschriebenen Gruppe \(\mathbb Z/3\mathbb Z\) sind die positiven
 Vielfachen von \(\overline1\) gleich
 \(\overline1,\overline2,\overline0,\overline1,\ldots\).
 In der Potenzschreibweise der Halbgruppe ist das
 \((\mu,\lambda)=(1,3)\). Die dritte Potenz ist das neutrale,
 also idempotente Element.
 \item[Ein Kreis mit einem Wert: \(\lambda=1\).]
 Die Potenzfolge wird ab \(\mu\) konstant. In der Nullhalbgruppe
 \(\{a,0\}\), in der jedes Produkt \(0\) ist und \(a\ne0\),
 lautet sie \(a,0,0,\ldots\).
 Hier ist \((\mu,\lambda)=(2,1)\).
 Die Operation ist assoziativ, weil jede Klammerung eines
 Dreifachprodukts \(0\) ergibt.
 \item[Schon der erste Wert ist idempotent.]
 Aus \(a^2=a\) folgt durch wiederholte Multiplikation
 \(a^r=a\) für alle positiven \(r\).
 Das Paar lautet \((1,1)\), und \(\langle a\rangle=\{a\}\).
\end{description}

Die beiden Parameter messen verschiedene Eigenschaften.
\(\mu\) sagt, wann der Kreis beginnt; \(\lambda\) sagt,
wie viele verschiedene Werte in diesem Kreis liegen.

\section*{Index und Periode aus einer Tafel ablesen}

Ist die assoziative Grundtafel gegeben, lässt sich das Paar ohne
Durchprobieren möglicher Parameter finden:
\begin{enumerate}
 \item Starte mit dem Wert \(a^1=a\) und notiere den Exponenten \(1\).
 \item Multipliziere den zuletzt erreichten Wert rechts mit \(a\).
       Dadurch erhältst du die nächste Potenz.
 \item Ist ihr Wert neu, notiere seinen Exponenten und setze fort.
 \item Taucht bei Exponent \(j\) zum ersten Mal ein schon
       notierter Wert auf, lies seinen früheren Exponenten \(i\) ab.
       Gib \(\mu=i\) und \(\lambda=j-i\) aus.
\end{enumerate}
Bei einem Träger mit \(N\) Elementen muss dieser letzte Schritt
spätestens beim Exponenten \(N+1\) eintreten.
Der erste Treffer liefert genau die im Existenzbeweis verwendete
Wiederholung; die vorherigen Werte waren verschieden.
Die Eindeutigkeit des Index--Perioden-Paars begründet, dass dieses
Verfahren die bestimmten Parameter des Elements findet.

Die Assoziativität der gegebenen Tafel ist hierbei Voraussetzung.
Für eine beliebige binäre Operation kann die Folge der wiederholten
Rechtsmultiplikationen ebenfalls einen Kreis bilden. Das allein
liefert noch kein Additionsgesetz für Potenzen und keine Halbgruppe.

\section*{Verbindungen zum Hauptband}

Der Beweisgang verwendet die Endlichkeit zuerst für die Existenz einer
Wiederholung. Sobald ein Index--Perioden-Paar gegeben ist, folgen die
periodische Fortsetzung und die Beschreibung der erzeugten Unterhalbgruppe
bereits aus den Halbgruppenaxiomen und den Eigenschaften dieses Paars.
Die formalen Aussagen im Hauptband halten diese Voraussetzungen getrennt.

Die \IndexPeriodProofsLink{} führen vom Nachweis der
Wohldefiniertheit über Existenz und Eindeutigkeit bis zu
Kardinalität und Idempotenz. Vor jeder Beweistabelle steht die
zu beweisende Aussage mit ihren Voraussetzungen.

Wie man anschließend die \emph{gesamte} endliche Halbgruppe aus
ihrer Tafel bis auf Umbenennung bestimmt, erklärt die
\TableIsomorphismsReadingLink{} zu Multiplikationstafeln und Isomorphie.
Für den größeren Zusammenhang mit Idealen und Potenzhalbgruppen steht
die \FiniteSemigroupsReadingLink{} zur Verfügung.

